\label{sec:item-difficulty-transfer}
We ask whether item features recover the ordering of observed item
difficulty. The target is the mean binary correctness per item, centered
within benchmark for fitting. We compare a ridge probe of item embeddings,
a probe using log character count, and a sampled random score. The ID
diagnostic holds out items within the 11 training benchmarks. The other
two columns evaluate entire unseen benchmarks. All methods use matched
items and sampled pairs. Pairwise accuracy excludes observed ties and
gives half credit for predicted ties. Results are averaged across benchmarks.

\begin{table}[t]
\centering
\caption{Item-difficulty transfer (WP7). Cells report within-benchmark
Spearman correlation and pairwise accuracy, respectively.}
\label{tab:item-difficulty}
\begin{tabular}{lccc}
\toprule
Feature & ID (11) & Validation (3) & Test (3)\\
\midrule
Embedding & $0.400$ / $0.649$ & $0.196$ / $0.570$ & $-0.020$ / $0.494$\\
Log character count & $0.107$ / $0.539$ & $0.058$ / $0.525$ & $0.061$ / $0.521$\\
Random score & $-0.003$ / $0.499$ & $-0.024$ / $0.490$ & $0.076$ / $0.529$\\
\bottomrule
\end{tabular}
\end{table}

Embeddings show a clear signal on known benchmarks. The historical figure
of about $0.09$ Spearman ($0.53$ pairwise) pools validation and test across
all six unseen benchmarks. On the three test benchmarks alone, neither
feature beats the matched sampled random control on either metric.
This is a descriptive result: the sampled control can fluctuate around
chance, and three benchmarks cannot establish that the features have no
transfer value.

